% Semantic Operating System. Place this master beside the three original inputs.
% Compile with Tectonic or XeLaTeX. The PDF includes each source as read at build time.
% The distributed sos.pdf is the compiled master under the requested publication name.
\documentclass[11pt,a4paper]{article}
\usepackage[T1]{fontenc}
\usepackage{lmodern}
\usepackage[a4paper,left=23mm,right=23mm,top=23mm,bottom=22mm,headheight=15pt,headsep=6mm,footskip=10mm]{geometry}
\usepackage{xcolor,booktabs,array,fancyhdr}
% Master preamble dependencies, alongside the existing page/header layout.
\usepackage{mathtools,amssymb,mathrsfs,needspace,docmute}
% input file is standalone; docmute skips its private preamble.
% XeTeX with legacy T1 fonts requires a TeX en-dash for this Unicode input.
\ifdefined\XeTeXversion
\begingroup
\lccode`\~="2013
\lowercase{\endgroup
  \newcommand{\OcaInputDash}{\catcode"2013=\active\def~{\textendash}}}
\fi
\newcommand{\RenderOriginalAlgebra}[1]{%
  \begingroup
  \ifdefined\XeTeXversion\OcaInputDash\fi
  \allowdisplaybreaks[2]%
  \setlength{\jot}{4pt}%
  \setlength{\parindent}{0pt}%
  \setlength{\parskip}{4pt}%
  \raggedright
  \renewcommand{\href}[2]{##2}%
  \setcounter{soslaw}{0}%
  \RenewDocumentCommand{\section}{sm}{\stepcounter{soslaw}\manualtarget{law-\arabic{soslaw}}\subsection*{##2}}%
  \input{#1}%
  \par\endgroup}
% Numbered source editions. Requires hyperref and xcolor.
\definecolor{SosListNumber}{HTML}{526770}
\definecolor{SosListFence}{HTML}{A6B9BE}
\ExplSyntaxOn
\ior_new:N \l__sos_source_ior
\int_new:N \l__sos_line_int
\bool_new:N \l__sos_fence_bool
\bool_new:N \l__sos_code_bool
\str_new:N \l__sos_key_str
\tl_new:N \l__sos_line_tl
\seq_new:N \l__sos_bold_seq
\seq_new:N \l__sos_italic_seq
\cs_generate_variant:Nn \seq_set_split_keep_spaces:Nnn { NnV }

% Legacy T1 fonts take typographic punctuation through their text commands.
\cs_new_protected:Npn \__sos_glyph:n #1
  {
    \int_case:nnF { `#1 }
      { {8211}{\textendash} {8217}{\textquoteright} }
      {#1}
  }

% Retain indentation and allow a long literal to wrap on the same numbered row.
\cs_new_protected:Npn \__sos_code:n #1
  {
    \str_map_inline:nn {#1}
      {
        \int_case:nnF { `##1 }
          {
            {9}{ \hskip 4\fontdimen2\font\relax }
            {32}{ \penalty 0\relax \hskip \fontdimen2\font\relax }
          }
          { \__sos_glyph:n {##1} \penalty 900\relax }
      }
  }

% Prose uses word spacing, with extra breakpoints in the original URLs.
\cs_new_protected:Npn \__sos_prose:n #1
  {
    \str_map_inline:nn {#1}
      {
        \__sos_glyph:n {##1}
        \int_case:nn { `##1 }
          {
            {47}{\penalty 100\relax}
            {45}{\penalty 100\relax}
            {95}{\penalty 100\relax}
            {46}{\penalty 100\relax}
          }
      }
  }
\cs_new_protected:Npn \__sos_emphasis:n #1
  {
    \seq_set_split_keep_spaces:Nnn \l__sos_italic_seq {*} {#1}
    \seq_map_indexed_inline:Nn \l__sos_italic_seq
      {
        \int_if_even:nTF {##1}
          { {\itshape \__sos_prose:n {##2}} }
          { \__sos_prose:n {##2} }
      }
  }
\cs_new_protected:Npn \__sos_markdown:n #1
  {
    \regex_match:nnTF { \A \x{23}+ \s } {#1}
      {
        \tl_set:Nn \l__sos_line_tl {#1}
        \regex_replace_once:nnN { \A \x{23}+ \s+ } {} \l__sos_line_tl
        \sffamily\bfseries \exp_args:NV \__sos_prose:n \l__sos_line_tl
      }
      {
        \tl_set:Nn \l__sos_line_tl {#1}
        % Print the cited authors and titles from Markdown links.
        \regex_replace_all:nnN { \[ ([^\]]+) \] \( https?:// [^)]+ \) } { \1 } \l__sos_line_tl
        \seq_set_split_keep_spaces:NnV \l__sos_bold_seq {**} \l__sos_line_tl
        \seq_map_indexed_inline:Nn \l__sos_bold_seq
          {
            \int_if_even:nTF {##1}
              { {\bfseries \__sos_emphasis:n {##2}} }
              { \__sos_emphasis:n {##2} }
          }
      }
  }
\cs_new_protected:Npn \__sos_linenumber:
  {
    \hbox to 6.5mm
      {
        \hfil
        \hypertarget{\str_use:N\l__sos_key_str-\int_use:N\l__sos_line_int}
          {
            {\sffamily\fontsize{7}{7}\selectfont\color{SosListNumber}
             \int_use:N\l__sos_line_int}
          }
      }
    \kern 1.3mm
    \bool_if:NTF \l__sos_fence_bool
      { {\color{SosListFence}\vrule width .45pt} }
      { \kern .45pt }
    \kern 1.7mm
  }
\cs_new_protected:Npn \__sos_source_row:n #1
  {
    \int_incr:N \l__sos_line_int
    \bool_set_true:N \l__sos_code_bool
    \str_if_eq:VnT \l__sos_key_str {M}
      {
        \str_if_eq:nnTF {#1} {```}
          { \bool_set_inverse:N \l__sos_fence_bool }
          { \bool_set_eq:NN \l__sos_code_bool \l__sos_fence_bool }
      }
    \par\nointerlineskip
    \hbox to \linewidth
      {
        \__sos_linenumber:
        \tl_if_blank:nTF {#1}
          { \vrule width 0pt height 5.5pt depth 1.5pt\hfil }
          {
            \vtop
              {
                \hsize=\dimexpr\linewidth-9.5mm-.45pt\relax
                \parindent=0pt\relax
                \raggedright\hyphenpenalty=200\relax
                \noindent
                \bool_if:NTF \l__sos_code_bool
                  {
                    \ttfamily\fontsize{8.4}{10.5}\selectfont
                    \strut \__sos_code:n {#1}\strut
                  }
                  {
                    \normalfont\fontsize{10}{12.3}\selectfont
                    \strut \__sos_markdown:n {#1}\strut
                  }
                \par
              }\hfil
          }
      }
    \penalty 0\relax
  }
\NewDocumentCommand{\OriginalListing}{mm}
  {
    \group_begin:
      \par\str_set:Nn \l__sos_key_str {#1}
      \int_zero:N \l__sos_line_int
      \bool_set_false:N \l__sos_fence_bool
      \ior_open:Nn \l__sos_source_ior {#2}
      \ior_str_map_inline:Nn \l__sos_source_ior { \__sos_source_row:n {##1} }
      \ior_close:N \l__sos_source_ior
      \iow_term:x { Source~#1:~\int_use:N\l__sos_line_int\space physical~lines }
    \group_end:
  }
\ExplSyntaxOff
\usepackage[unicode,bookmarksnumbered=false,colorlinks=true,linkcolor=tealref,urlcolor=tealref]{hyperref}
\definecolor{tealref}{HTML}{155B61}
\definecolor{ink}{HTML}{192C37}
\definecolor{quiet}{HTML}{465D66}
\definecolor{rulegray}{HTML}{B8C8C8}
\hypersetup{pdftitle={Semantic Operating System},pdfsubject={Structure-substrate names, Object Counting Algebra and observer-dependent descriptions}}
\setlength{\parindent}{0pt}
\setlength{\parskip}{4.5pt plus 1pt}
\setlength{\emergencystretch}{2em}
\raggedbottom
\newcommand{\EditionPart}{Formal discussion}
\pagestyle{fancy}\fancyhf{}
\fancyhead[L]{\sffamily\small\color{quiet}Semantic Operating System}
\fancyhead[R]{\sffamily\small\color{quiet}\EditionPart}
\fancyfoot[L]{\sffamily\footnotesize\color{quiet}\ifdefined\SosManualEdition\hyperlink{start}{Start}\quad\hyperlink{dictionary}{Dictionary}\quad\hyperlink{index}{Index}\else\hyperlink{guide}{Reading guide}\fi}
\fancyfoot[R]{\sffamily\small\color{ink}\thepage}
\renewcommand{\headrulewidth}{0.3pt}
\newcommand{\kicker}[1]{{\sffamily\bfseries\small\color{tealref}#1\par}}
\newcommand{\edtitle}[1]{{\sffamily\bfseries\fontsize{24}{27}\selectfont\color{ink}#1\par}}
\newcommand{\spread}[2]{\clearpage\renewcommand{\EditionPart}{#1}\kicker{Formal discussion \enspace/\enspace #1}\edtitle{#2}\vspace{2mm}}
\newcounter{soslaw}
\newcommand{\manualtarget}[1]{\hypertarget{#1}{}\label{loc-#1}}
\newcommand{\manualpage}[1]{\hyperlink{#1}{\pageref*{loc-#1}}}
\newcommand{\appendixpart}[3]{\clearpage\renewcommand{\EditionPart}{#2}\pdfbookmark[1]{#2}{bm#1}\manualtarget{#1}\kicker{#3}}
\newcommand{\soschapter}[4]{\clearpage\renewcommand{\EditionPart}{#2}\pdfbookmark[1]{#4}{bm#1}\manualtarget{#1}\kicker{#3}\edtitle{#4}\smallskip}
\newcommand{\sosterm}[3]{\par\Needspace{2\baselineskip}\manualtarget{#1}\textbf{#2.}\enspace #3\par}
\newcounter{question}
\newcommand{\question}[2]{\refstepcounter{question}\par\addvspace{4mm}\manualtarget{q\thequestion}\pdfbookmark[2]{\thequestion. #1}{bmq\thequestion}%
  {\sffamily\bfseries\color{tealref}Question \thequestion\enspace\textcolor{rulegray}{/}\enspace #1\par}%
  {\large\bfseries\color{ink}#2\par}\nobreak}
\newcommand{\ans}{\noindent{\sffamily\bfseries\small\color{tealref}Answer \thequestion.}\enspace\ignorespaces}
\newcommand{\linesref}[3]{\hyperlink{#1-#2}{#1\,#2--#3}}
\newcommand{\lineref}[2]{\hyperlink{#1-#2}{#1\,#2}}
\newcommand{\sources}[1]{\par\nobreak{\raggedright\sffamily\footnotesize\color{quiet}Source lines\enspace #1\par}}
\newcommand{\pairlink}[2]{\hyperlink{q#1}{#2}}
\newcommand{\N}{\mathbb N_0}
\newcommand{\Cspace}{\mathcal C}
\newcommand{\Name}{\mathsf{Name}}
\newcommand{\Boolean}{\mathbb B}
\begin{document}
\color{ink}
\ifdefined\SosManualEdition
\hypersetup{pdftitle={Semantic Operating System: Operational Manual, v2},pdfsubject={An operational guide with algorithms, terminology, equations and line-numbered sources}}
\manualtarget{start}\pdfbookmark[0]{Semantic Operating System: Operational Manual}{bmstart}
\renewcommand{\EditionPart}{Operational manual}
\kicker{Operational manual \enspace/\enspace Version 2}
\edtitle{Semantic Operating System}
{\large How to store a named structure, read it from an observer's point of origin, and calculate with its cells\par}
\medskip
\textbf{A first run.} Use two primitive types, neutral and toggle, in that order. Give a record the name \texttt{lamp}, shape $(1,1)$ and cell value $(1,1)$. Those last two numbers say the cell contains one occurrence of each type. Shape and content happen to have the same numbers here. They have different roles.

\begin{tabular}{@{}p{.29\linewidth}p{.63\linewidth}@{}}
\toprule
\textbf{Action} & \textbf{Expected observation}\\
\midrule
Register the record & The name \texttt{lamp} resolves to one occupied position.\\
Call with completion $0$ & The returned pair is $(\texttt{lamp},1)$. The second field asks for more detail.\\
Read the sole cell & A detailed observer sees $(1,1)$. A size observer sees $2$. The cell's natural canonical address is $4$.\\
Divide the cell by itself & Both singleton types reconstruct it. The exact quotient addresses are $\{1,2\}$.\\
\bottomrule
\end{tabular}

\textbf{Follow the operation.} The formal discussion explains why these values differ. Five implementations give the steps for reproducing them with other names, shapes and algebra profiles. Click a topic or use the printed page number.
\begin{center}\small
\begin{tabular}{@{}p{.65\linewidth}r@{}}
\toprule
\textbf{Where to go} & \textbf{Page}\\
\midrule
\hyperlink{guide}{Twelve questions and answers} & \manualpage{guide}\\
\hyperlink{implementations}{Implementations and worked traces} & \manualpage{implementations}\\
\hyperlink{dictionary}{Key terminology and notation} & \manualpage{dictionary}\\
\hyperlink{equations}{The fourteen sections of Object Counting Algebra} & \manualpage{equations}\\
\hyperlink{draft}{Operating-system draft, numbered source D} & \manualpage{draft}\\
\hyperlink{discussed}{oca-discussion, numbered source M} & \manualpage{discussed}\\
\hyperlink{texlisting}{OCA formulation, numbered source T} & \manualpage{texlisting}\\
\hyperlink{index}{Alphabetical subject index} & \manualpage{index}\\
\bottomrule
\end{tabular}
\end{center}
\clearpage\renewcommand{\EditionPart}{Formal discussion}
\fi
\manualtarget{guide}\pdfbookmark[1]{Reading guide and formal discussion}{bmguide}
\kicker{Formal specification \enspace/\enspace 22 November 2026}
\ifdefined\SosManualEdition\edtitle{The formal discussion}\else\edtitle{Semantic Operating System}\fi
{\large A discussion of named structures, observer views and Object Counting Algebra\par}
\vspace{2mm}
\textbf{Reading guide.} The twelve answers define a mathematical interface for the draft. A choice introduced to make that interface usable is identified where it first matters. The same two primitive objects, \emph{neutral} and \emph{toggle}, recur from the first record to the final call.
\par{\small\sffamily
\hyperlink{draft}{D\enspace Operating-system draft}\quad
\hyperlink{discussed}{M\enspace oca-discussion}\quad
\hyperlink{texlisting}{T\enspace OCA formulation}\\
Each source citation opens its numbered line. The \hyperlink{equations}{typeset algebra} preserves the fourteen original mathematical sections.\par}
\smallskip\hrule height .3pt
\question{The record}{What does a structure-substrate name identify?}
\ans It identifies a record whose detail an observer can query. To give the draft's $n$ dimensions a precise meaning, take a record $R=(s,\lambda,h,B)$. Here $s$ is its exact text name, $\lambda$ is a fixed text key for the algebra used by its cells, and $h=(h_1,\ldots,h_n)$ gives its extents, with $n\geq1$ and $h_j\in\N$. The body $B$ assigns one configuration in $\Cspace_\lambda$ to each coordinate in $\prod_{j=1}^{n}\{0,\ldots,h_j-1\}$. A zero extent gives an empty body with a known shape. A valid coordinate may instead contain the empty configuration. For example, the record named \texttt{lamp} can have shape $(1,1)$ and contain one neutral and one toggle in its sole cell. The counts, the cell coordinate and the record name answer three different lookup questions.
\sources{\lineref{D}{2}, \lineref{D}{5}\quad\linesref{T}{73}{120}\quad\linesref{M}{318}{336}}
\question{Reversible text}{How can the name or the entire record travel as one integer?}
\ans Fix text to finite sequences of Unicode scalar values, with exact equality. This edition adopts UTF-8 and preserves spelling, order and whitespace. For a text $s$, prefix its bytes with byte \texttt{01} and read the result as an unsigned big-endian integer $E(s)$. To invert, write the integer in its shortest big-endian form, check and remove the prefix, and decode valid UTF-8. Reject values outside the image. The empty text maps to $1$, and a single zero byte maps to $256$. The names \texttt{ab} and \texttt{ba} receive different codes even though counting their letters would give the same multiset.

For a whole record, frame each text field by its UTF-8 byte length in decimal, a hash sign, then the field itself. Concatenate frames for $s,\lambda,n$, the $n$ extents, then each cell's multiplicities in alphabet order, visiting cells in row-major order (last axis fastest). Use decimal integers with no leading zeros except $0$ itself. A negative coefficient in the signed profile uses an ASCII minus followed by its positive magnitude. Reject signed zero. Apply $E$ to the resulting text. A decoder uses the fixed profile $\lambda$ to determine the count of fields per cell, verifies lengths and shape, and rejects missing or trailing fields. This defines a reversible record code even when a name contains a hash sign.
\sources{\lineref{D}{2}\quad\linesref{M}{662}{678}, \linesref{M}{726}{746}}
\spread{Configuration and rank}{Counting a cell}
\question{The counting domain}{Once the record is found, what does a cell's canonical number identify?}
\ans Fix the finite, nonempty alphabet $A=(o_1,\ldots,o_d)$ in namespace $\lambda$. Its natural configuration space is $\Cspace=\N^d$. Coordinate $x_i$ counts occurrences of $o_i$, the zero vector is the empty multiset, and $\nu(x)=\sum_i x_i$ counts all occurrences. Two configurations are equal exactly when every corresponding count agrees. A record can move in the registry without changing any of these values.

Order the entire space by increasing $\nu$ and then by ascending lexicographic order on $(x_1,\ldots,x_d)$. The rank $\rho(x)$ counts the configurations preceding $x$. Each cardinality level is finite, so the ranks start at zero, are consecutive, and have an inverse $\delta$. With neutral first and toggle second, the first six addresses are
\[
\begin{array}{c|rrrrrr}
a&0&1&2&3&4&5\\ \hline
\delta(a)&(0,0)&(0,1)&(1,0)&(0,2)&(1,1)&(2,0).
\end{array}
\]
The \texttt{lamp} cell has two occurrences and address $4$, even if it is the first entry registered. An OCA address requires its namespace. The integer encoding of \texttt{lamp} in Answer~2 requires its text codec.
\sources{\linesref{T}{73}{120}, \linesref{T}{125}{185}, \linesref{T}{889}{916}\quad\linesref{M}{680}{704}, \linesref{M}{714}{746}}
\question{Aggregation and removal}{Given two cell addresses, when can the system remove what it added?}
\ans Decode both addresses using $\delta$. Aggregation adds corresponding counts, $(x\oplus y)_k=x_k+y_k$. It is total, associative and commutative, has zero as its identity, and admits cancellation. Exact removal $x\ominus y$ is defined precisely when $y_k\leq x_k$ for every $k$. Its result is the unique $z$ for which $z\oplus y=x$. Encode a successful result with $\rho$.

Adding addresses $1$ and $2$ therefore gives address $4$. Removing the toggle at address $1$ from $4$ recovers address $2$. Attempting to remove that toggle from address $2$ fails because $(1,0)$ contains zero toggles. Removing a toggle from itself succeeds with the empty vector, address $0$. These are operations on specified cells. They leave the record name, its shape and the registry untouched. Further storage effects would need a separate write operation.
\sources{\linesref{T}{190}{242}, \linesref{T}{247}{278}, \linesref{T}{433}{476}\quad\linesref{M}{362}{404}, \linesref{M}{568}{574}}
\spread{Composition and division}{Reconstructing a configuration}
\question{Composition}{How does the system know what two primitive objects produce together?}
\ans The algebra profile supplies $T_{ij}^{k}\in\N$, the number of occurrences of type $k$ produced by the ordered primitive pair $(i,j)$. For each such pair, $x_i y_j$ counts its interactions. Summing their contributions defines the output
\[
                 (x\star y)_k=\sum_{i=1}^{d}\sum_{j=1}^{d}x_i y_j T_{ij}^{k}.
\]
This distributes over aggregation on both sides and makes zero absorbing. Associativity requires agreement on the two groupings of every primitive triple. Commutativity requires $T_{ij}^{k}=T_{ji}^{k}$ for every $i,j,k$. A two-sided unit must reproduce each basis configuration on both sides. If each pair produces one primitive and the table forms a finite monoid with identity, its identity singleton is the unit, cardinalities multiply, and the configuration algebra is the nonnegative monoid semiring.

In the running instance neutral acts as identity and two toggles compose to neutral. The product of two address-$1$ operands is therefore address $2$. A profile that leaves even one ordered primitive pair unspecified does not determine this operation on all cells.
\sources{\linesref{T}{190}{214}, \linesref{T}{247}{300}, \linesref{T}{305}{372}, \linesref{T}{673}{748}\quad\linesref{M}{406}{450}, \linesref{M}{560}{574}}
\question{Exact division}{If the product and one factor are known, what must an answer contain?}
\ans Every admissible factor that reconstructs the product. With divisor $y$ on the right, $Q_{\rm R}(x,y)=\{q:q\star y=x\}$. With $y$ on the left, $Q_{\rm L}(x,y)=\{q:y\star q=x\}$. For each side the condition is a linear system $M(y)q=x$ over nonnegative integers. A consistent system with a zero column has infinitely many solutions because its corresponding coordinate is free. Otherwise a nonempty fiber is finite because each coordinate contributes a positive amount to a fixed total.

Return the encoded set as $\mathsf{Sol}(\rho[Q])$. For \texttt{lamp}, dividing $(1,1)$ by $(1,1)$ on either side requires $q_1+q_2=1$. Neutral and toggle both work, so the result is $\mathsf{Sol}(\{1,2\})$. A nonzero dividend divided by zero gives $\mathsf{Sol}(\varnothing)$, while zero divided by zero gives $\mathsf{Sol}(\N)$ in the natural namespace. A symbolic set can denote the latter exactly. The evaluator distinguishes a scalar $\mathsf{Val}(0)$, $\mathsf{Sol}(\{0\})$, $\mathsf{Sol}(\varnothing)$ and a failed residual $\mathsf{Fail}$.
\sources{\linesref{T}{231}{242}, \linesref{T}{377}{428}, \linesref{T}{481}{534}, \linesref{T}{917}{934}\quad\linesref{M}{452}{480}, \linesref{M}{574}{580}}
\spread{Interpretation and namespaces}{Changing what an observer sees}
\question{Observer origin}{What detail may differ between observers of the same named substrate?}
\ans Give each origin $o$ a declared view $V_o$. In this edition it retains the record's name and shape and applies a map $\sigma_o:\Cspace_\lambda\to S_o$ to each cell. A request for \texttt{lamp} can return the pair $(1,1)$ to a type-counting observer and the integer $2$ to an observer who sees only cardinality. Both read the same cell. An observation has its own representation and need not let the observer recover the complete record. Reversible serialization in Answer~2 applies to the chosen record or text.

To compute algebraically on meanings, specify $0_o,\oplus_o,\star_o$ and require $\sigma_o(0)=0_o$, $\sigma_o(x\oplus y)=\sigma_o(x)\oplus_o\sigma_o(y)$ and $\sigma_o(x\star y)=\sigma_o(x)\star_o\sigma_o(y)$. Then every exact factor maps to a semantic factor on the same side, though a semantic solution may lack a corresponding exact factor. Cardinality makes singleton neutral and singleton toggle indistinguishable. Removing a toggle from the latter succeeds, while removing it from the former fails. Consequently an observer who retains only total size cannot decide exact removal from that view.
\sources{\lineref{D}{2}\quad\linesref{T}{539}{599}, \linesref{T}{601}{635}\quad\linesref{M}{506}{534}, \linesref{M}{748}{762}}
\question{Changes of domain}{What survives relabeling, restriction, or the admission of negative counts?}
\ans A bijection $\pi$ on the address space relabels the same configuration by $\pi(\rho(x))$. Transport the decoder, operations, quotient sets and interpretation through the same bijection and each reconstruction still holds. Keep the namespace with each stored address.

Restricting to $D\subseteq\Cspace$ containing zero changes which configurations exist for that profile. Rank within $D$, accept a scalar operation only if its result remains in $D$, and intersect each quotient set with $D$. For $D=\{(0,0),(1,1)\}$ the mixed cell has restricted rank $1$ and adding it to itself has no admitted result. In the full namespace the same cell has rank $4$.

The signed extension has domain $\mathbb Z^d$, total coordinatewise subtraction and its own order by $\ell^1$ norm followed by lexicographic order. In that profile the mixed-cell division of Answer~6 admits every $(t,1-t)$ for $t\in\mathbb Z$. The nonnegative profile admits exactly two of them. A caller must select the domain before interpreting the addresses or comparing answers.
\sources{\linesref{T}{640}{668}, \linesref{T}{753}{810}, \linesref{T}{815}{884}, \linesref{T}{922}{934}\quad\linesref{M}{536}{556}, \linesref{M}{726}{746}}
\spread{Lookup and completion}{Reading the operating-system draft}
\question{Registry lookup}{How does a name reach an array cell without becoming its array index?}
\ans The draft asks for an $n$-dimensional store indexed by name. Supply a finite registry grid $G$ with declared extents, a partial injective index $L:\Name\rightharpoonup G$, and a partial array $P:G\rightharpoonup\mathsf{Record}$. Require $\operatorname{im}L=\operatorname{dom}P$ and $P(L(s)).\mathrm{name}=s$ for every stored name. Together these conditions give each registered name one occupied position and every occupied position one name. $P$ stores records, each with its own inner shape from Answer~1. The depth of that inner array and the rank of the registry grid are independent parameters in this edition.

For an exact, nonempty name $s$, lookup reads $P(L(s))$ if $s\in\operatorname{dom}L$ and reports \textsf{UnknownName} otherwise. The empty text has a codec value but is inadmissible as a registry name in this profile. A duplicate name or occupied coordinate is invalid when preparing the registry. Movement to another free coordinate preserves lookup if the index and array change together. All queries below use one fixed, valid snapshot. The draft specifies neither a registration algorithm nor a persistence or concurrency protocol.
\sources{\linesref{D}{2}{9}\quad\linesref{M}{662}{678}, \linesref{M}{714}{746}}
\question{Completion and learning}{What does the written status function return for the draft's two Boolean cases?}
\ans The first test accepts $1$ and assigns $0$ to \texttt{learn}. The second test repeats the same condition after the first one has failed, so its assignment of $1$ is unreachable. With input $0$ the return statement supplies no defined value for \texttt{learn}. For a total interface this edition adopts the complementary second case. Interpret \texttt{learn} as a request for further detail and \texttt{complete} as the caller's declaration.
\begin{center}\small
\begin{tabular}{@{}lll@{}}
\toprule
Input \texttt{complete} & Literal draft result & Adopted \texttt{status\_ping} result\\
\midrule
$0$ (incomplete) & no value specified & $1$ (request further detail)\\
$1$ (complete) & $0$ & $0$ (no request)\\
\bottomrule
\end{tabular}
\end{center}
For Boolean input $c\in\{0,1\}$ the adopted rule is $\operatorname{status\_ping}(c)=1-c$. Out-of-domain input is rejected at the interface. Calling this function computes a flag. It does not inspect a record, infer its completeness from the observer, or start a learning process.
\sources{\linesref{D}{13}{27}, \linesref{D}{29}{33}}
\spread{Interface and trace}{Following a call to its result}
\question{The return contract}{How do lookup and status compose in \texttt{structure\_substrate}?}
\ans The draft places a name and a status call in its return expression. This edition makes those two fields an ordered result. For a well-formed name $s$ and Boolean $c$ in a valid snapshot, resolve $s$ with Answer~9. On success return $(s,\operatorname{status\_ping}(c))$. On an absent name return \textsf{UnknownName}. Reject malformed text, an empty name or a non-Boolean flag as \textsf{InvalidInput} before lookup. A successful call preserves both the stored body and the caller's original name. Asking for observer-specific detail is a separate read of the resolved record through $V_o$.

The placeholders, \texttt{elif}, and undeclared \texttt{learn} in the draft are pseudocode. In C++ a built-in comma expression evaluates to its right operand, so a direct translation of line~9 would lose the name. A typed pair or a record with fields \texttt{name} and \texttt{learn} implements the declared result. Its text field owns its value in this interface. The profile specifies a mathematical read with no side effects or storage writes.
\sources{\linesref{D}{5}{11}, \linesref{D}{13}{25}\quad\linesref{M}{662}{678}}
\question{An end-to-end trace}{What happens when the draft's entry point asks about \texttt{lamp}?}
\ans Line~31 supplies input $0$. Bind its name placeholder to \texttt{lamp} and prepopulate a registry position with the one-cell record of Answer~1. In the literal draft the function reaches its return with an undefined status. Under the adopted contract, lookup finds \texttt{lamp}, the ping returns $1$, and the variable \texttt{x1} receives $(\texttt{lamp},1)$. The \texttt{return 0} in \texttt{main} is a separate process result. The draft contains no print operation.

A second call on the same snapshot with input $1$ yields $(\texttt{lamp},0)$. Both calls leave the cell at natural address $4$. A subsequent detailed observation recovers $(1,1)$, its size view returns $2$, and dividing that cell by itself returns the two addresses $\{1,2\}$. Removing the name from the example snapshot would instead yield \textsf{UnknownName} at lookup. These predictions distinguish lookup, learning status, observer detail and algebraic results well enough to check each independently.
\sources{\linesref{D}{29}{35}, \linesref{D}{13}{25}\quad\linesref{T}{889}{934}\quad\linesref{M}{558}{580}}
\ifdefined\SosManualEdition
% Each procedure reads one fixed namespace and uses exact, unbounded integer arithmetic.
\newcommand{\sosstep}[1]{\par\addvspace{3pt}\noindent{\sffamily\bfseries\color{tealref}#1}\enspace\ignorespaces}
\newcommand{\sosop}[1]{\textsf{#1}}
\soschapter{implementations}{Implementations}{Implementation 1 of 5}{Reversible record codec}
\manualtarget{impl-1}\label{sos-impl-one}
\begingroup\small\raggedright
The five procedures implement one SOS, from record input to observed output. The last procedure combines them into three complete operating profiles. The codec first converts exact text or a complete record into an integer and recovers the input. Supply a fixed catalog of algebra profiles keyed by exact text. A profile declares an ordered alphabet of length $d$, its count domain and any admissibility test. This catalog is the shared input to all five implementations.

\sosstep{Encode and decode text.}
For each Unicode scalar, write its nonnegative integer value in binary. Choose the row for its value, left-pad with zero bits to fill the slots marked $b$, and insert the bits from most to least significant. Concatenate the bytes in order. Hexadecimal bounds below use digits 0--9 and A--F for values 10--15.
\begin{center}\small
\begin{tabular}{@{}ll@{}}\toprule
Scalar value (hexadecimal) & UTF-8 byte patterns, in order\\\midrule
\texttt{0--7F} & \texttt{0bbbbbbb}\\
\texttt{80--7FF} & \texttt{110bbbbb 10bbbbbb}\\
\texttt{800--FFFF}, excluding \texttt{D800--DFFF} &
\texttt{1110bbbb 10bbbbbb 10bbbbbb}\\
\texttt{10000--10FFFF} & \texttt{11110bbb 10bbbbbb 10bbbbbb 10bbbbbb}\\\bottomrule
\end{tabular}\end{center}
To decode, read the bytes from first to last. Match each leading byte, consume its continuation bytes, join their $b$ bits and check that the value belongs to the selected row. Reject truncation, invalid bytes and values outside the row. To obtain $E(s)$, start $v=1$ and update $v\gets256v+b$ for each byte $b$. To invert $E$, require a positive integer. Repeatedly take quotient and remainder by 256 until the quotient is zero. Read the remainders in reverse order, require first byte 1, and UTF-8-decode the remaining bytes.

\sosstep{Pack a record.}
Validate the nonempty name $s$, known profile $\lambda$, dimension $n\geq1$, nonnegative extents $h_1,\ldots,h_n$, and one admissible $d$-vector per cell. Traverse cells with the last coordinate changing fastest. Form text fields in the order $s,\lambda,n,h_1,\ldots,h_n$, followed by each vector's coordinates in alphabet order. For each field $f$ append its UTF-8 byte length in decimal, the character \texttt{\#}, then $f$. Return $E$ of the concatenated text. Integers use ASCII digits without leading zeros except \texttt{0}. A negative signed count uses a minus followed by a nonzero magnitude.

\sosstep{Unpack an integer.}
Invert $E$ and parse the resulting bytes from the beginning. Each frame consumes a canonical decimal length, a \texttt{\#}, and exactly that many valid UTF-8 bytes. Read the first three fields, resolve $\lambda$, read $n$ extents, then $d\prod_j h_j$ counts. Require each field's declared type, the record checks above, and exact end of input. Return the reconstructed record or \sosop{InvalidInput}. A zero extent consumes zero cell fields.

\sosstep{Checks.}
$E(\text{empty text})=1$. A single zero byte gives 256. Encoding an empty registry name fails. A \texttt{\#} inside a field survives decoding because the length fixes its end. For every accepted record $R$, unpacking its packed integer returns $R$. Canonical framing also makes repacking an accepted integer reproduce it.
\sources{\linesref{D}{2}{5}\quad\linesref{M}{662}{678}\quad\hyperlink{q2}{Answer 2}}
\endgroup

\soschapter{impl-2}{Implementations}{Implementation 2 of 5}{Named registry and status}
\label{sos-impl-two}
\begingroup\small\raggedright
The registry accepts records from \hyperlink{impl-1}{the codec}. It returns a new snapshot when a registration succeeds. Passing the chosen snapshot explicitly makes repeated reads reproducible.

\sosstep{State.}
A snapshot $S=(G,L,P)$ has a finite coordinate grid $G$, a partial map $L$ from exact nonempty names to positions, and a partial array $P$ from positions to records. Grid indices start at zero. An empty registry has $L=P=\varnothing$. In a supplied snapshot, verify that $L$ is injective, that its occupied positions equal the keys of $P$, that all these positions lie in $G$, and that every record $P(L(s))$ validates and has name $s$. Reject a malformed snapshot before serving requests. A record body has its own independent shape.

\sosstep{Register$(S,g,v)$.}
Here $g$ is a grid coordinate and $v$ is a packed record. Apply the following steps in order.
\begin{enumerate}\setlength{\itemsep}{1pt}
\item Unpack $v$ with implementation 1. Require $g$ to have the grid's number of indices and each index to be an integer in bounds. On either failure return \sosop{InvalidInput}.
\item Let $R$ be the decoded record. If $R.s$ is present in $L$, return \sosop{DuplicateName}. If $g$ is present in $P$, return \sosop{OccupiedPosition}.
\item Construct $L'=L\cup\{R.s\mapsto g\}$ and $P'=P\cup\{g\mapsto R\}$. Return $(G,L',P')$ together, as one new snapshot. A rejected operation returns its error and leaves $S$ available to the caller.
\end{enumerate}

\sosstep{Lookup$(S,s)$.}
Require exact, valid, nonempty text. Reject malformed names as \sosop{InvalidInput}. If $s$ is absent from $L$, return \sosop{UnknownName}. Otherwise return $P(L(s))$. Equality compares the sequence of scalar values exactly, including case and whitespace.

\sosstep{StructureSubstrate$(S,s,c)$.}
Validate the name and the Boolean $c\in\{0,1\}$ before looking up $s$. Pass through a lookup error. After success return the pair $(s,1-c)$ with its own copy of the name. The second field is the adopted total \texttt{status\_ping} result. The original draft repeats its test for $c=1$ and leaves the $c=0$ path unspecified.

\sosstep{Checks.}
Register a valid record named \texttt{lamp} at an unoccupied $g$. Lookup in the old empty snapshot returns \sosop{UnknownName}. Lookup in the returned snapshot retrieves the original record. Registering the same name again returns \sosop{DuplicateName}. In the returned snapshot, flags $0$ and $1$ yield $(\texttt{lamp},1)$ and $(\texttt{lamp},0)$. Flag $2$ yields \sosop{InvalidInput} even for an absent name.
\sources{\linesref{D}{5}{11}, \linesref{D}{13}{33}\quad\hyperlink{q9}{Answer 9}\quad\hyperlink{q11}{Answer 11}}
\endgroup

\soschapter{impl-3}{Implementations}{Implementation 3 of 5}{Rank and evaluate a cell}
\label{sos-impl-three}
\begingroup\small\raggedright
The input is a fixed namespace $\lambda$, an operation, and two addresses tagged with $\lambda$. The profile declares dimension $d\geq1$, the tensor $T_{ij}^{k}\in\mathbb N_0$ for every input pair and output index, and a natural or signed count domain. A restriction supplies a terminating membership test and either its finite domain size or a declaration that the domain is infinite. A relabeling supplies mutually inverse permutations $\pi,\pi^{-1}$ of the rank space. Use identity functions when labels are canonical.

\sosstep{Enumerate configurations.}
Visit shells $t=0,1,\ldots$. In a natural shell, visit tuples in $\{0,\ldots,t\}^d$ in lexicographic order and retain those with $\sum_i x_i=t$. In a signed shell, visit $\{-t,\ldots,t\}^d$ in the same order and retain those with $\sum_i |x_i|=t$. In either case emit only tuples admitted by the profile. Lexicographic traversal increments the last coordinate fastest. Every admissible tuple occurs once, at a finite position. The loops prioritize simplicity over speed.

\sosstep{Rank$(\lambda,x)$ and Unrank$(\lambda,a)$.}
For rank, require an admitted vector with $d$ integer components. Count emitted tuples preceding $x$, starting the counter at zero, and return $\pi$ of this count. For unrank, require a matching namespace and a nonnegative integer address. When the declared size is $m<\infty$, reject $a\ge m$ immediately. Apply $\pi^{-1}$ to $a$ and return the tuple emitted at that zero-based position. A bad address, mismatched namespace, or malformed request returns \sosop{InvalidInput}. These contracts give $\operatorname{Unrank}(\lambda,\operatorname{Rank}(\lambda,x))=x$ and the converse for valid addresses.

\sosstep{Evaluate$(\lambda,\omega,a,b)$.}
Validate the operation, unrank both inputs as $x,y$, and propagate input errors. For addition or subtraction, form $z_k=x_k+y_k$ or $z_k=x_k-y_k$. For multiplication, initialize all $z_k$ to zero and for every $i,j,k$ add $x_i y_jT_{ij}^{k}$ to $z_k$. For these three operations, return $\mathsf{Fail}$ if the resulting vector falls outside the profile, including a negative natural residual. Otherwise return $\mathsf{Val}(\operatorname{Rank}(\lambda,z))$. For right or left division dispatch $(\lambda,x,y,\text{side})$ to \hyperlink{impl-4}{implementation 4}.

\sosstep{Checks.}
With the natural neutral--toggle table of \linesref{T}{889}{914}, the mixed vector $(1,1)$ has rank 4. Adding addresses 1 and 2 gives $\mathsf{Val}(4)$, and removing address 1 from address 4 gives $\mathsf{Val}(2)$. Removing address 1 from 2 gives $\mathsf{Fail}$. Restrict the carrier to $\{(0,0),(1,1)\}$ to obtain rank 1 for the same mixed vector and \sosop{InvalidInput} when decoding address 2.
\sources{\linesref{T}{125}{185}, \linesref{T}{190}{242}, \linesref{T}{481}{534}, \linesref{T}{753}{810}, \linesref{T}{851}{871}}
\endgroup

\soschapter{impl-4}{Implementations}{Implementation 4 of 5}{Reconstruct every factor}
\label{sos-impl-four}
\begingroup\small\raggedright
Given a side and the validated vectors $x,y$ from \hyperlink{impl-3}{the evaluator}, return $\mathsf{Sol}$ of the ranks of \emph{all} admissible $q$ solving $Aq=x$. For right division set $A_{ki}=\sum_j T_{ij}^k y_j$. For left division set $A_{ki}=\sum_j T_{ji}^k y_j$. An infinite answer is represented by a parameterized set.

\sosstep{Natural counts.}
For column $i$ compute $c_i=\sum_k A_{ki}$. Write $F=\{i:c_i=0\}$ for the free coordinates. Nonnegativity makes each such column identically zero. For every other $i$ enumerate $0\le q_i\le\lfloor(\sum_k x_k)/c_i\rfloor$ in a finite Cartesian product. Fill the free entries with zeros and retain a base vector only if $Aq=x$. The result contains every retained base with its free entries independently replaced by arbitrary nonnegative integers. When there are no bounded coordinates, test the one empty tuple. With no retained bases the answer is empty. With any base and any free coordinate it is infinite in the unrestricted natural profile. The bounds follow by summing the rows of $Aq=x$.

\sosstep{Signed counts.}
Use the following recursion for an arbitrary integer matrix $A$ with $m$ columns and target $b$ (initially $x$).
\begin{enumerate}\setlength{\itemsep}{0pt}\setlength{\topsep}{2pt}
\item Drop every all-zero row with zero target. An all-zero row with nonzero target returns the empty set. With no rows left, return $\mathbb Z^m$ including the single empty tuple when $m=0$.
\item Start $U$ as the $m$-square identity. Pick a remaining row, move it and its target entry to the top, and swap a nonzero column into position one. Apply that column swap and every subsequent column operation to both $A$ and $U$.
\item While another entry $a_j$ of the first row is nonzero, put $a=A_{11}$. Choose the integer $k$ with $0\le a_j-ka<|a|$. Subtract $k$ times column one from column $j$. If its new top entry is nonzero, swap columns $j$ and one. Each such swap decreases the pivot magnitude.
\item Only $a=A_{11}\ne0$ now remains in the top row. Return empty unless $a$ divides $b_1$. Fix $p=b_1/a$ and recurse on the lower rows and columns $2,\ldots,m$, with target $b_i-A_{i1}p$ in each lower row. Return $U(p,t)$ for every tuple $t$ the recursion returns. The reversible column operations preserve every integer solution in these coordinates.
\end{enumerate}

\sosstep{Apply the namespace and check.}
Intersect the resulting family with the profile's admissible set and take its pointwise rank image. For a finite restriction one may enumerate its members and test $Aq=x$ directly. An infinite restriction can retain membership as a symbolic filter. When $y=0$, the ambient result is the whole carrier for $x=0$ and empty otherwise. In the neutral--toggle example, dividing $(1,1)$ by itself gives $q_1+q_2=1$. The natural identifier set is $\{1,2\}$. The signed configurations are $(t,1-t)$ for all integers $t$.
\sources{\linesref{T}{377}{428}, \linesref{T}{773}{777}, \linesref{T}{843}{849}, \linesref{T}{922}{934}}
\endgroup

\soschapter{impl-5}{Implementations}{Implementation 5 of 5}{Dispatch an observer query}
\label{sos-impl-five}
\begingroup\small\raggedright
Use a fixed snapshot $S$, the profile catalog, and an observer table. An observer entry $(o,\lambda)$ supplies a total map $\sigma_o$ on admitted configurations. The request specifies an exact name $s$ and one of the following operations. The response carries the operation's result type.

\sosstep{Dispatch$(S,s,\text{request})$.}
Reject an unrecognized operation, a malformed name or wrongly typed arguments as \sosop{InvalidInput}. Process a status request with \hyperlink{impl-2}{StructureSubstrate}. A full-record request returns \hyperlink{impl-1}{Pack} of the looked-up record. Other operations start with \hyperlink{impl-2}{Lookup} and propagate \sosop{UnknownName} for an absent, well-formed name.
\begin{enumerate}\setlength{\itemsep}{1pt}
\item For $\sosop{View}(o)$, resolve the registered map for $(o,R.\lambda)$ or return \sosop{InvalidInput}. Visit the cells of $R$ in row-major order. In a newly allocated output array of the same shape place $\sigma_o(B[g])$ at each $g$. Return the original name and shape with that output array. A zero extent yields an empty view of that shape.
\item For $\sosop{Cell}(g)$, require an integer coordinate for each inner axis with $0\le g_j<h_j$. Otherwise return \sosop{InvalidInput}. Return the address of $B[g]$ from \hyperlink{impl-3}{Rank}, tagged with $R.\lambda$.
\item For $\sosop{Calculate}(g,\omega,b)$, resolve the cell address and pass it as the first operand to \hyperlink{impl-3}{Evaluate} with the supplied second address $b$. Both addresses use the cell's namespace. Division returns \hyperlink{impl-4}{a solution set}. These reads use the stored multiplicities, including when an earlier view combined distinguishable configurations.
\end{enumerate}

\sosstep{Operating profile A. Natural inventory.}
Use namespace \texttt{nt}, all nonnegative count pairs, the canonical rank and alphabet (neutral, toggle). Neutral is the identity and two toggles compose to neutral. Pack the record \texttt{lamp} of shape $(1,1)$ with sole cell $(1,1)$. The exact text to which $E$ is applied is
\begin{center}\ttfamily 4\#lamp2\#nt1\#21\#11\#11\#11\#1\end{center}
Register the integer in an empty one-position registry at $(0)$. In the resulting snapshot, a status request with $c=0$ returns $(\texttt{lamp},1)$. At inner coordinate $(0,0)$, a cell request returns $(\texttt{nt},4)$. A view using $\sigma(x)=\sum_i x_i$ returns \texttt{lamp}, shape $(1,1)$ and sole value 2. A calculation at $(0,0)$ requesting right division by $(\texttt{nt},4)$ returns $\mathsf{Sol}(\{1,2\})$ in \texttt{nt}. Either factor, multiplied on the right by the original mixed cell, reconstructs it. Unpacking a full-record response recovers the original record.

\sosstep{Operating profile B. Signed balance.}
Give namespace \texttt{bal} one primitive $u$ with $u\star u=u$, carrier $\mathbb Z$, and the canonical signed rank. Pack a one-cell record \texttt{net} holding $-2$, register it at $(0)$ in a fresh grid, and observe the signed count itself. Its cell has address 3. Subtract address 2 (the count $+1$) to obtain $\mathsf{Val}(5)$, which decodes to $-3$. A status request with $c=1$ returns $(\texttt{net},0)$.

\sosstep{Operating profile C. Capacity register.}
For capacity $K\ge1$ choose a dedicated namespace, the same one-primitive table and carrier $D=\{0,\ldots,K\}$ in ascending rank order. Pack and register a one-cell record named \texttt{stock} holding $K$ at $(0)$ in a fresh grid. Observe remaining capacity as $K-x$. For $K=2$ the view gives 0. Adding address 1 to the cell returns $\mathsf{Fail}$, and dividing the cell by address 1 returns $\mathsf{Sol}(\{2\})$.

\sosstep{Failure check for profile A.}
Coordinate $(0,1)$ is outside the inner shape and yields \sosop{InvalidInput}. Running the same well-formed name request on the old empty registry yields \sosop{UnknownName}. These interface errors occur before algebraic evaluation.
\sources{\linesref{D}{29}{35}\quad\linesref{T}{539}{599}, \linesref{T}{889}{934}\quad\hyperlink{q12}{Answer 12}}
\endgroup

% Included by the SOS operating manual. All links resolve inside the manual.
\soschapter{dictionary}{Terminology dictionary}{Reference}{Key terminology and notation}
A name, a cell coordinate and a configuration address answer different lookup questions. The entries below explain each term in the operational interface and the symbols used in the original equations.

\noindent\textsf{\hyperlink{g-a}{A--C}\quad\hyperlink{g-d}{D--H}\quad
\hyperlink{g-i}{I--N}\quad\hyperlink{g-o}{O--Q}\quad
\hyperlink{g-r}{R--S}\quad\hyperlink{g-t}{T--Z}\quad
\hyperlink{g-notation}{Mathematical signs}}

\newcommand{\sosdictgroup}[2]{\par\Needspace{5\baselineskip}\manualtarget{#1}\subsection*{#2}}

\sosdictgroup{g-a}{A--C}
\sosterm{g-addition}{Addition, $\oplus$}{Combine two configurations by adding matching coordinates. With $x=(1,0)$ and $y=(0,1)$, $x\oplus y=(1,1)$. Zero is the identity. Addition admits cancellation, meaning $x\oplus z=y\oplus z$ implies $x=y$.}
\sosterm{g-admissible}{Admissible}{Allowed by the selected profile. In the natural profile every coordinate is a nonnegative integer. A restricted profile can impose additional conditions on whole configurations. An operation on admitted inputs may still produce an excluded result. See \hyperlink{g-restricted}{restricted domain}.}
\sosterm{g-lattice}{Affine lattice}{A set of integer vectors consisting of a particular solution plus every integer combination of some direction vectors. For signed neutral--toggle division, $(0,1)+t(1,-1)$ for $t\in\mathbb Z$ describes an entire quotient set.}
\sosterm{g-alphabet}{Alphabet, $A$}{The finite, nonempty ordered list $(o_1,\ldots,o_d)$ of primitive types. Its order fixes which type a coordinate counts and breaks ties in the canonical ranking. The size $d$ counts types, including types whose current multiplicity is zero.}
\sosterm{g-array}{Array, axis and cell}{An axis is one coordinate position in an array index. A cell is the value at one permitted index tuple. The registry array holds records. The array inside a record holds configurations. See \hyperlink{g-shape}{shape and extent}.}
\sosterm{g-associative}{Associativity and commutativity}{Associativity means $(x\star y)\star z=x\star(y\star z)$. The predicate $\mathsf A_T$ tests equality for every triple of basis configurations and every output coordinate. Commutativity means $x\star y=y\star x$. The predicate $\mathsf C_T$ tests $T_{ij}^{k}=T_{ji}^{k}$ for every $i,j,k$. Bilinearity extends these basis tests to all configurations. \linesref{T}{305}{346}.}
\sosterm{g-atomic}{Atomic registration}{Publishing the name index and the corresponding occupied grid cell as one state change. A reader of a snapshot can see either complete registry state. A failed registration leaves the preceding state available.}
\sosterm{g-basis}{Basis configuration, $e_i$}{The vector containing one occurrence at coordinate $i$ and zero elsewhere. Every natural configuration has the unique decomposition $\bigoplus_i x_i e_i$. The expression $n x$ means the sum of $n$ copies of $x$, with zero copies giving the empty configuration.}
\sosterm{g-bijection}{Bijection}{A function pairing each input with one output so that every output has exactly one input. Its inverse reverses the pairing. \emph{Injective} means distinct inputs have distinct outputs. \emph{Surjective} means every member of the declared output domain is reached.}
\sosterm{g-bilinearity}{Bilinearity and distributivity}{Composition distributes over a sum in either operand. Thus every term $x_i y_j$ contributes $x_i y_j$ copies of $e_i\star e_j$ to the product. Both products with zero equal zero, the absorbing element for multiplication.}
\sosterm{g-byte}{Bit, byte and big-endian order}{A bit is zero or one. A byte has eight bits and represents a number from $0$ to $255$. In big-endian order the first byte gives the coefficient of the greatest power of $256$. For example, the bytes \texttt{01 00}, shown in hexadecimal (base sixteen), represent $256$.}
\sosterm{g-rank}{Canonical rank and unranking, $\rho,\delta$}{The rank $\rho(x)$ counts configurations before $x$ in the declared order. The natural profile orders by size and then lexicographically. The inverse $\delta=\rho^{-1}$ recovers the vector. In the rank formula $p_i(x)=\sum_{h<i}x_h$ counts earlier coordinates. To unrank an address, find its size shell and select the vector at its position within that shell. \linesref{T}{125}{185}.}
\sosterm{g-cardinality}{Cardinality, $\nu$}{For a natural configuration $x$, $\nu(x)=\sum_i x_i$ counts occurrences. For a set $U$, $|U|$ counts members. Thus $|\{(1,1)\}|=1$ while $\nu((1,1))=2$. On addresses, $\nu_I=\nu\circ\delta$. Product sizes multiply for all $x,y$ exactly when every primitive pair yields one occurrence.}
\sosterm{g-codec}{Codec, $E$}{A matched encoder and decoder. For SOS text, $E(s)$ is the unsigned integer formed from the single prefix byte \texttt{01} followed by the UTF-8 bytes of $s$. Decode the shortest big-endian representation, check its prefix, then decode the remaining bytes. The empty text has code $1$. Record encoding applies $E$ after framing its fields. See \hyperlink{impl-1}{Implementation~1}.}
\sosterm{g-completion}{Completion and learning flag}{The caller supplies a Boolean $c$, either $0$ or $1$. The adopted $\operatorname{status\_ping}(c)=1-c$ yields $1$ to request further detail and $0$ when complete. Draft lines 15 and 19 both test $1$, leaving $c=0$ unresolved in the original. \linesref{D}{13}{25}.}
\sosterm{g-tensor}{Composition tensor, $T$, and product, $\star$}{An entry $T_{ij}^{k}$ counts occurrences of type $k$ produced by one ordered primitive pair $(o_i,o_j)$. The vector $\mu_T(o_i,o_j)$ has these entries as its coordinates. For configurations, $(x\star y)_k=\sum_{i,j}x_i y_j T_{ij}^{k}$. The entire tensor belongs to the profile.}
\sosterm{g-configuration}{Configuration, $\mathcal C$}{A multiset of typed occurrences represented by a multiplicity vector. In the natural case $\mathcal C=\mathbb N_0^d$. Two vectors are equal exactly when corresponding counts agree. The zero vector $\mathbf 0$ has no occurrences. The signed extension uses integer coefficients.}
\sosterm{g-congruence}{Congruence and semantic equivalence}{Write $x\sim_\sigma y$ when $\sigma(x)=\sigma(y)$. This equivalence is a congruence for an operation when replacing arguments by equivalent ones gives equivalent results. Defining a partial operation on equivalence classes also requires compatible definedness. Cardinality can hide whether a particular type can be removed. \linesref{T}{601}{635}.}
\sosterm{g-order}{Coordinatewise order, $\preccurlyeq$}{The relation $y\preccurlyeq x$ means $y_i\leq x_i$ for every coordinate $i$. It checks availability of every requested type for exact removal. For example, neither $(1,0)$ nor $(0,1)$ lies below the other.}

\sosdictgroup{g-d}{D--H}
\sosterm{g-determinant}{Determinant, $\det M$}{A signed number associated with a square matrix. For a $1$-by-$1$ matrix take its only entry. For larger matrices expand along row one, summing $(-1)^{1+j}M_{1j}$ times the determinant after deleting row one and column $j$. A nonzero determinant makes $Mv=0$ have only the zero rational solution. For a fixed natural divisor and side this is equivalent to at most one quotient for every dividend. A zero determinant permits a collision for some dividend. \linesref{T}{423}{428}.}
\sosterm{g-diagonalization}{Diophantine elimination and pivots}{A Diophantine equation asks for integer solutions. In the signed solver, a pivot is the nonzero entry used to simplify a row. The algorithm swaps columns and subtracts integer multiples of one column from another, preserving an invertible integer coordinate change (a unimodular map). When the row has a single nonzero coefficient, that pivot must divide the corresponding target. Fix its coordinate and solve the remaining rows recursively. See \hyperlink{impl-4}{Implementation~4}.}
\sosterm{g-division}{Division}{Exact multiplicative reconstruction. Declare which side holds the known factor and return all admitted solutions of that equation. See \hyperlink{g-quotient}{quotient set} and \hyperlink{g-matrix}{multiplication matrix}.}
\sosterm{g-domain}{Domain, image and closure}{The domain $\operatorname{dom} f$ contains inputs where $f$ has a value. The image $\operatorname{im} f$ consists of its attained values. A total function is defined for every allowed input. A partial one may be undefined for some. A set is closed under an operation when applying it to members always yields another member.}
\sosterm{g-errors}{Errors and empty results}{\textsf{InvalidInput} rejects malformed arguments, invalid encodings, or numbers outside the declared input domain. \textsf{UnknownName} means a valid name is absent. On registration, \textsf{DuplicateName} means the name already exists and \textsf{OccupiedPosition} means the requested position is taken. Within the algebra, $\mathsf{Fail}$ signifies an undefined scalar operation. $\mathsf{Sol}(\varnothing)$ reports a valid division with no solutions.}
\sosterm{g-evaluator}{Evaluator and outcome tags}{The set of operation symbols $\Omega$ contains $\Omega_{\rm alg}=\{\oplus,\star,\ominus\}$ and the two divisions. $\operatorname{Eval}(\omega,a,b)$ returns a tagged outcome in $\mathsf{Out}_I$. $\mathsf{Val}(a)$ carries one address, $\mathsf{Sol}(U)$ carries a set of addresses, and $\mathsf{Fail}$ records an undefined scalar result. Thus $\mathsf{Val}(0)$ and $\mathsf{Sol}(\{0\})$ have distinct types. \linesref{T}{481}{534}.}
\sosterm{g-fiber}{Fiber}{The complete collection of inputs sent to a given output by a function. Fixing $y$ gives a function $q\mapsto q\star y$. Its fiber over $x$ is $Q_{\rm R}(x,y)$. A fiber can be empty, finite or infinite.}
\sosterm{g-monoid}{Finite monoid and monoid algebra}{A monoid is a set with associative multiplication and a two-sided identity. In $\mathsf M_{m,h}(T)$ the rule $m(i,j)$ chooses the unique output type of each primitive pair, is associative, and has identity $h$. Hence $T_{ij}^{k}=\mathbf1_{\{m(i,j)=k\}}$. Write $M=([d],m,h)$. Its monoid algebra $\mathbb N_0[M]$ uses natural coefficients and the product induced by $m$. The algebra $\mathbb Z[M]$ permits signed coefficients. The unit is $e_h$. \linesref{T}{673}{723}.}
\sosterm{g-framing}{Framing and row-major order}{One field is its canonical decimal UTF-8 byte length, then \texttt{\#}, then exactly that many bytes. For example, \texttt{3\#a\#b} encodes the field \texttt{a\#b}. Record fields are name, profile, axis count, extents and cell coefficients. Visit cells in row-major order, varying the last axis fastest. The shape and alphabet give the number of coefficient fields.}
\sosterm{g-free}{Free commutative monoid, $\operatorname{FreeCMon}(A)$}{The multiset domain under addition and zero. Every map from the alphabet to a commutative monoid extends in exactly one addition- and identity-preserving way to the configurations. Thus counts completely determine each sum of primitives.}
\sosterm{g-free-coordinate}{Free coordinate}{A component of a solution vector that can vary independently within its declared number domain. A zero column in a nonnegative multiplication matrix gives a free natural coordinate whenever a solution exists.}
\sosterm{g-gr}{Group completion, $\operatorname{Gr}$}{Pairs of natural configurations represent formal differences. Pairs $(x,y)$ and $(x',y')$ represent the same difference exactly when $x\oplus y'=x'\oplus y$. Taking one class for each equivalence gives $\mathbb Z^d$ under addition, a group because every vector has an additive inverse. \linesref{T}{815}{841}.}
\sosterm{g-homomorphism}{Homomorphism}{A map preserving the stated operations and identities. For the semantic predicate $\mathsf H_\sigma$, see \hyperlink{g-semantics}{semantic interpretation}. A unital homomorphism also preserves the designated multiplicative identity.}

\sosdictgroup{g-i}{I--N}
\sosterm{g-address}{Identifier and sequential indexing, $I$}{An identifier is the address of one configuration within a namespace. Sequential identifiers use successive labels. In the natural profile $I=\mathbb N_0$ and the canonical order determines which configuration receives each label. A registry assigning the next free label on arrival can be sequential while its labels depend on history.}
\sosterm{g-identity}{Identity, zero and unit}{An identity leaves the other operand unchanged on both sides. The additive identity is $\mathbf0$, with address $0_I=\rho(\mathbf0)$. $\mathcal U_T$ is the set of two-sided multiplicative identities. Checking $u\star e_i=e_i\star u=e_i$ for every $i$ suffices. There is at most one $u$. If it exists its address is $1_I=\rho(u)$, which is $2$ in the neutral--toggle example.}
\sosterm{g-kernel}{Kernel}{For a matrix $M$ on a specified coefficient domain, its zero fiber is the set of $v$ satisfying $Mv=0$. Over the integers, adding any such $v$ to one solution of $Mq=x$ gives another. For a meaning map the associated equivalence groups configurations with equal images.}
\sosterm{g-lex}{Lexicographic order, $<_{\rm lex}$}{Compare two coordinate vectors at their first differing position, from left to right. The vector with the smaller coordinate there comes first. Natural OCA uses this comparison only after sizes tie, giving every configuration finitely many predecessors. The total order is denoted $\prec$.}
\sosterm{g-matrix}{Multiplication matrix, $M_\varepsilon(y)$}{For a fixed divisor $y$, column $i$ of $M_{\rm R}(y)$ is $e_i\star y$. Column $i$ of $M_{\rm L}(y)$ is $y\star e_i$. With either choice, $M_\varepsilon(y)q=x$ is the quotient equation. The column sum $c_i^\varepsilon(y)$ counts output occurrences contributed by one unit at $i$. Every natural solution satisfies $\sum_i c_i^\varepsilon(y)q_i=\nu(x)$. \linesref{T}{377}{421}.}
\sosterm{g-multiset}{Multiset, multiplicity and occurrence}{An occurrence is one instance of a primitive type. Its type's multiplicity is the number of instances of that type. A multiset records each multiplicity, so changing the order in which occurrences are listed preserves configuration identity.}
\sosterm{g-name}{Name}{The exact original text of a structure-substrate. Registered names are nonempty. Name equality compares the sequence of Unicode scalar values, including every space, in order. Looking up a name and encoding that name as an integer are independent operations.}
\sosterm{g-namespace}{Namespace and profile, $\lambda$}{A namespace fixes the configuration domain, its equality criterion and its indexing convention, including alphabet order. A profile also supplies the composition rules and any admissibility constraints. The record stores the profile's text key $\lambda$. A bare address can denote different vectors in different namespaces.}
\sosterm{g-neutral-toggle}{Neutral and toggle}{Two primitive types in the running instance, ordered neutral first. Neutral is the identity, and toggle composed with toggle gives neutral. For vectors $(p,q),(r,s)$ the product is $(pr+qs,ps+qr)$. The mixed vector $(1,1)$ has natural rank $4$. \linesref{T}{889}{934}.}

\sosdictgroup{g-o}{O--Q}
\sosterm{g-oca}{Object Counting Algebra, $\mathfrak O(A,T,\mathscr S,\sigma)$}{The specification collecting configurations and their operations, address encoding and decoding, the transported operations, semantic interpretation and the typed evaluator. The parameters state which primitive types exist, what their pairs produce, and how results are interpreted. \linesref{T}{48}{66}.}
\sosterm{g-observer}{Observer origin and view, $o,V_o$}{An origin names a declared way of observing a record. Its view keeps the name and shape and applies a meaning map $\sigma_o$ to each cell. A type-counting view can return $(1,1)$ where a size-only view returns $2$. The record survives either observation.}
\sosterm{g-operation}{Operation and operand}{An operation is a rule taking inputs, called operands, to a result when defined. The symbols $x,y,z,q$ usually denote configurations and $a,b$ their integer addresses. An \emph{ordered} pair remembers which operand is first. The side of a divisor matters when multiplication can be noncommutative.}
\sosterm{g-primitive}{Primitive type and primitive rule}{A type $o_i$ is one element of the alphabet. Its singleton configuration is $e_i$. The rule $\mu_T$ assigns a whole output configuration to each ordered pair of types. It may produce zero, one or many occurrences. The finite monoid specialization always produces one.}
\sosterm{g-quotient}{Quotient set, $Q_{\rm R},Q_{\rm L}$}{A quotient is an admissible missing factor. $Q_{\rm R}(x,y)$ collects all $q$ with $q\star y=x$, whereas $Q_{\rm L}(x,y)$ collects all $q$ with $y\star q=x$. The subscript locates the \emph{known} divisor $y$. For natural mixed neutral--toggle divided by itself, both encoded sets equal $\{1,2\}$. A symbolic family denotes a set by a rule generating exactly its members.}

\sosdictgroup{g-r}{R--S}
\sosterm{g-record}{Record or structure-substrate, $R$}{The tuple $R=(s,\lambda,h,B)$ holds an exact name, profile key, shape and body. The body maps every valid coordinate to a configuration. A record with a zero extent has no cells. A record with one cell containing $\mathbf0$ has that cell and its empty configuration.}
\sosterm{g-registry}{Registry, $G,L,P$}{The outer array and its name index. The grid $G$ contains the permitted registry positions. The partial injective map $L$ sends a name to an occupied position. The partial array $P$ stores the record there. The conditions $\operatorname{im}L=\operatorname{dom}P$ and $P(L(s)).\mathrm{name}=s$ give exactly one name per occupied position. See \hyperlink{impl-2}{Implementation~2}.}
\sosterm{g-relabel}{Relabeling, $\pi$}{A bijection on the address space changes the encoder to $\rho_\pi=\pi\circ\rho$ and the decoder to $\delta_\pi=\delta\circ\pi^{-1}$. Transport every operation, constant, solution set and interpretation with this change. Then corresponding operands have corresponding results. The symbol $\cong_\pi$ denotes this isomorphism. \linesref{T}{640}{668}.}
\sosterm{g-residual}{Residual, $x\ominus y$ and $\mathcal R(x,y)$}{The exact additive remainder is the unique vector $z$ with $z\oplus y=x$, if one exists. In the natural profile it exists exactly when $y\preccurlyeq x$, and then $z_i=x_i-y_i$. The residual \emph{set} $\mathcal R(x,y)$ contains that vector when defined and is empty otherwise.}
\sosterm{g-restricted}{Restricted domain, $D$}{A declared subset of $\mathcal C$ containing zero. For $\omega_D$, inputs and result must belong to $D$ and the ambient operation must be defined. $Q_{\varepsilon,D}=Q_\varepsilon\cap D$ on admitted operands. The rank $\rho_D$ counts predecessors in $D$ alone. $I_D=\rho_D[D]$ and $\delta_D=\rho_D^{-1}$. If $|D|=k$ is finite, its addresses are $0,\ldots,k-1$. \linesref{T}{753}{810}.}
\sosterm{g-ring}{Ring}{A semiring in which every element has an additive inverse. For the signed algebra $\mathsf A_T$ and existence of a two-sided unit suffice to give a unital ring, written $\mathsf{Ring}_1$ in T. Commutativity of multiplication is a separate property.}
\sosterm{g-semantics}{Semantic interpretation and compatibility}{The target $\mathscr S=(S,\oplus_S,\star_S,0_S)$ gives possible meanings, their operations and designated zero. The map $\sigma:\mathcal C\to S$ assigns meanings. $\sigma_I=\sigma\circ\delta$ interprets addresses. Predicate $\mathsf H_\sigma$ requires preservation of zero, sums and products. $\mathcal R_S,Q_{\rm R}^{S},Q_{\rm L}^{S}$ solve the corresponding reconstruction equations in $S$. Under $\mathsf H_\sigma$ every interpreted configuration solution is a semantic solution (soundness). The reverse inclusion would require every such semantic solution to be the image of an admitted configuration solution (completeness). \linesref{T}{539}{599}.}
\sosterm{g-semiring}{Semiring}{A unital semiring has associative, commutative addition with zero, associative multiplication with a unit, distributivity on both sides, and multiplicatively absorbing zero. $\mathsf{Semiring}_1$ denotes this class. $\mathsf{CSemiring}_1$ adds commutativity of multiplication. \linesref{T}{348}{366}.}
\sosterm{g-shape}{Shape, extent and coordinate}{The shape $h=(h_1,\ldots,h_n)$ gives the number of positions along each of the $n$ axes. A coordinate $k$ satisfies $0\le k_j<h_j$ on every axis. There are $\prod_j h_j$ cells. In shape $(2,3)$ the last cell in row-major order has coordinate $(1,2)$. Here $n$ counts axes, while $d$ counts primitive types in each cell.}
\sosterm{g-shell}{Shell and combination count}{The natural size shell $\mathcal C_n$ contains the vectors of size $n$. Its size $\binom{n+d-1}{d-1}$ counts arrangements of $n$ identical occurrences separated into $d$ coordinates by $d-1$ dividers. For $0\le k\le m$, $\binom{m}{k}$ counts selections of $k$ positions among $m$. T uses zero when $m<k$ and an empty sum equals zero.}
\sosterm{g-signed}{Signed configuration, $\mathcal C^\pm$}{An integer vector $z\in\mathbb Z^d$ has possibly negative coefficients. Subtraction is total. Composition uses the same bilinear formula and division admits all integer solutions. $\|z\|_1=\sum_i|z_i|$ is its absolute size. The order $\prec_\pm$ uses this size and then lexicographic order to give the new encoder and decoder $\rho_\pm,\delta_\pm$.}
\sosterm{g-snapshot}{Snapshot}{A fixed valid registry state used throughout one read or trace. An immutable snapshot is held constant after publication. A later registration can publish a new snapshot. A caller using the old one continues to see its original name index and body values.}
\sosterm{g-sos}{Structure-substrate query, \texttt{structure\_substrate}}{The adopted read checks the arguments, resolves the name in the snapshot, then returns the pair containing the same name and its completion-derived learning flag. Reading a view or evaluating an algebra operation is a separate request. The original draft's entry point supplies completion $0$.}

\sosdictgroup{g-t}{T--Z}
\sosterm{g-transport}{Transport and hatted notation}{A hat marks an operation on addresses. For example, $a\widehat\star b=\rho(\delta(a)\star\delta(b))$. A quotient set is encoded member by member, $\widehat Q_\varepsilon(a,b)=\rho[Q_\varepsilon(\delta(a),\delta(b))]$. The rule works for empty and infinite sets as well. \linesref{T}{433}{476}.}
\sosterm{g-utf8}{Unicode scalar and UTF-8}{A scalar is an integer from hexadecimal \texttt{0000} through \texttt{10FFFF} excluding the surrogate range \texttt{D800}--\texttt{DFFF}. UTF-8 uses the shortest suitable byte pattern. The first byte has form \texttt{0xxxxxxx}, \texttt{110xxxxx}, \texttt{1110xxxx}, or \texttt{11110xxx} for lengths one through four. The remaining bytes have form \texttt{10xxxxxx}. Put the scalar's bits into the $x$ positions from most to least significant, left-padding with zeros. A decoder checks length, continuation bytes, shortest form and scalar range.}
\sosterm{g-view}{View}{The result of applying a particular observer's declared projection to a record. Views encode which detail is visible. Recovering the original cell from a projected value requires that its projection distinguish the relevant configurations. See \hyperlink{g-observer}{observer origin}.}
\sosterm{g-zero-column}{Zero column}{A column of $M_\varepsilon(y)$ whose entries all vanish. Over natural coefficients this is equivalent to $c_i^\varepsilon(y)=0$. If a solution exists, its coordinate $i$ can then be increased arbitrarily, giving countably many solutions. With no zero column, each $q_i$ is bounded above by $\lfloor\nu(x)/c_i^\varepsilon(y)\rfloor$.}

\manualtarget{g-notation}\subsection*{Mathematical signs}
\begin{description}
\item[Numbers and indices.] $\mathbb N_0=\{0,1,2,\ldots\}$ and $\mathbb Z$ includes all integers. $[d]=\{1,\ldots,d\}$. The letters $i,j,k,\ell,m$ can index coordinate positions. $\mathbf1_{\{i=k\}}$ is $1$ when the condition holds, $0$ otherwise. $\lfloor t\rfloor$ is the greatest integer at most $t$. The symbol $r\mid s$ means $s=rk$ for some integer $k$, and $r\nmid s$ means no such integer exists.
\item[Sets.] $\in$ means membership, $\subseteq$ inclusion, $\cap$ intersection, $\varnothing$ the empty set, and $\sqcup$ a disjoint union. $\mathcal P(U)$ contains all subsets of $U$. $\{q:F(q)\}$ denotes every $q$ satisfying condition $F$. $\aleph_0$ is the cardinality of $\mathbb N_0$. A set of this size can be put in bijection with the nonnegative integers.
\item[Maps and products.] $f:X\to Y$ gives one value in $Y$ for each member of $X$. The arrow $\rightharpoonup$ allows partiality. $f[U]=\{f(u):u\in U\}$. $g\circ f$ applies $f$ first. $\operatorname{id}_X$ returns its input. $X^d$ is the set of $d$-tuples over $X$. $\prod_j X_j$ is the set of tuples choosing one member from each $X_j$. $\sum$ adds terms and $\prod$ multiplies numeric factors.
\item[Assertions.] $:=$ introduces a definition, $\forall$ means for every, $\exists$ means there exists, and $\exists!$ means there is exactly one. $\land$ joins simultaneous requirements, $\lor$ allows either, $\Rightarrow$ denotes implication, and $\iff$ denotes implication in both directions. $\cong$ denotes an invertible correspondence preserving the stated operations or structure.
\end{description}


\fi
\appendixpart{equations}{Typeset algebra}{Original formulation \enspace/\enspace T}
\RenderOriginalAlgebra{oca-defined.tex}
\appendixpart{draft}{Operating-system draft}{Numbered original \enspace/\enspace D}
\edtitle{Operating-system draft}
{\sffamily\small\texttt{semantic-operating-system-draft.txt}\par}\smallskip
\OriginalListing{D}{semantic-operating-system-draft.txt}
\appendixpart{discussed}{oca-discussion}{Numbered original \enspace/\enspace M}
\edtitle{oca-discussion}
{\sffamily\small\texttt{oca-discussed.md}\par}\smallskip
\OriginalListing{M}{oca-discussed.md}
\appendixpart{texlisting}{OCA formulation source}{Numbered original \enspace/\enspace T}
\edtitle{Object Counting Algebra formulation}
{\sffamily\small\texttt{oca-defined.tex}\par}\smallskip
\OriginalListing{T}{oca-defined.tex}
\ifdefined\SosManualEdition\soschapter{index}{Subject index}{Reference}{Subject index}
The middle column takes you to a definition or a formal result. The last column takes you to a procedure, discussion or worked example. All page numbers refer to this edition.

\newcommand{\idxrow}[3]{#1 & \manualpage{#2} & \manualpage{#3}\\}
\begingroup\small\renewcommand{\arraystretch}{1.10}
\begin{tabular}{@{}p{.62\linewidth}rr@{}}
\toprule
\textbf{Subject} & \textbf{Definition} & \textbf{Use}\\
\midrule
\idxrow{Address, canonical}{g-rank}{impl-3}
\idxrow{Addition and aggregation}{law-3}{q4}
\idxrow{Admissible domain and closure}{g-restricted}{q8}
\idxrow{Alphabet and alphabet order}{law-1}{q3}
\idxrow{Array shape, extents and empty cells}{g-shape}{q1}
\idxrow{Associativity and commutativity}{law-5}{q5}
\idxrow{Boolean completion and learning flag}{g-completion}{impl-2}
\idxrow{Canonical indexing and unranking}{law-2}{impl-3}
\idxrow{Cardinality, total number of occurrences}{law-1}{q7}
\idxrow{Cell configuration and equality}{g-configuration}{q3}
\idxrow{Codec, byte order and round trip}{g-codec}{impl-1}
\idxrow{Composition and primitive output tensor}{g-tensor}{q5}
\idxrow{Decoder errors and absent names}{g-errors}{impl-2}
\idxrow{Division, exact left and right quotients}{g-quotient}{impl-4}
\idxrow{Empty result, empty configuration and failure}{g-evaluator}{q6}
\idxrow{Fiber, finite or infinite solution set}{g-fiber}{impl-4}
\idxrow{Framing a record}{g-framing}{impl-1}
\idxrow{Identifier arithmetic and transport}{g-transport}{impl-3}
\idxrow{Interpretation and observer equivalence}{g-observer}{q7}
\idxrow{Matrix and free quotient coordinates}{law-6}{impl-4}
\idxrow{Monoid and multiplicative identity}{law-11}{q5}
\idxrow{Name, registry and snapshot}{g-registry}{impl-2}
\idxrow{Namespace and rank stability}{g-namespace}{q8}
\idxrow{Neutral and toggle}{g-neutral-toggle}{q12}
\idxrow{Observer query and view}{g-observer}{impl-5}
\idxrow{Reconstruction from a known factor}{g-quotient}{q6}
\idxrow{Record and structure-substrate}{g-record}{q1}
\idxrow{Residual and exact subtraction}{g-residual}{q4}
\idxrow{Restricted configurations}{g-restricted}{impl-5}
\idxrow{Semantic compatibility}{g-semantics}{law-9}
\idxrow{Signed configurations}{g-signed}{impl-4}
\idxrow{Status ping and the draft's repeated guard}{g-completion}{q10}
\idxrow{Typed evaluation outcomes}{g-evaluator}{impl-5}
\idxrow{UTF-8 and Unicode scalar values}{g-utf8}{impl-1}
\idxrow{Zero as dividend or divisor}{law-6}{q6}
\bottomrule
\end{tabular}
\endgroup

\soschapter{formula-index}{Procedure and formula index}{Reference}{Procedure and formula index}
\textbf{Operating procedures.} Read these in order to follow a record from its serialized form through a complete observer query.

\newcommand{\soslocator}[2]{\hyperlink{#1}{#2} & \manualpage{#1}\\}
\begingroup\small\renewcommand{\arraystretch}{1.13}
\begin{tabular}{@{}p{.76\linewidth}r@{}}
\toprule\textbf{Procedure} & \textbf{Page}\\\midrule
\soslocator{impl-1}{Encode and recover a record}
\soslocator{impl-2}{Register a name, look it up and return its status}
\soslocator{impl-3}{Convert a cell to an address and evaluate an operation}
\soslocator{impl-4}{Recover every factor consistent with a product}
\soslocator{impl-5}{Dispatch a query and inspect worked profiles}
\bottomrule
\end{tabular}

\medskip
\textbf{Original algebra.} Each topic opens the typeset mathematics. The line range opens the corresponding source in the numbered formulation.

\newcommand{\soslawrow}[4]{\hyperlink{law-#1}{#2} & \manualpage{law-#1} & \linesref{T}{#3}{#4}\\}
\begin{tabular}{@{}p{.58\linewidth}rr@{}}
\toprule\textbf{Section} & \textbf{Page} & \textbf{Source lines}\\\midrule
\soslawrow{1}{Primitive domain}{73}{120}
\soslawrow{2}{Canonical sequential indexing}{125}{185}
\soslawrow{3}{Primitive composition and configuration operations}{190}{242}
\soslawrow{4}{Unconditional algebraic laws}{247}{300}
\soslawrow{5}{Associativity, commutativity and identity}{305}{372}
\soslawrow{6}{Division as nonnegative integer fibers}{377}{428}
\soslawrow{7}{Transport to identifier arithmetic}{433}{476}
\soslawrow{8}{Typed outcomes and total evaluation}{481}{534}
\soslawrow{9}{Semantic interpretation and compatibility}{539}{635}
\soslawrow{10}{Reindexing invariance}{640}{668}
\soslawrow{11}{Finite monoid specialization}{673}{748}
\soslawrow{12}{Admissibility restriction}{753}{810}
\soslawrow{13}{Signed extension}{815}{884}
\soslawrow{14}{Neutral--toggle instance}{889}{934}
\bottomrule
\end{tabular}
\endgroup
\fi
\end{document}
